% ============================================================
%  PHỤ LỤC A: CHỨNG MINH CÁC ĐỊNH LÝ TRONG KHÓA LUẬN
% ============================================================

\chapter{Chứng minh các kết quả ở Mục 3.3}
\label{appx:proofs-misclustering}

Phụ lục này trình bày chứng minh chi tiết của các bổ đề và mệnh đề được sử dụng trong phân tích tỉ lệ phân cụm sai ở Mục 3.3 của Chương~\ref{chap:thuat-toan}. Các phát biểu được giữ nguyên ở phần chính, người đọc có thể quay lại đây để xem chi tiết kỹ thuật.

\subsection*{Chứng minh Bổ đề~\ref{lem:spectrum-CBCt} (Phổ của $CBC^\top$)}

\begin{proof}
Phổ của $B$ rút ra trực tiếp từ tác động lên $\mathbf{1}_K$ và lên phần bù trực giao.

Với $v = \mathbf{1}_K$,
\[
    B \mathbf{1}_K = \alpha \mathbf{1}_K + \beta \mathbf{1}_K (\mathbf{1}_K^\top \mathbf{1}_K) = (\alpha + K \beta) \mathbf{1}_K.
\]

Với mọi $v \in \mathbb{R}^K$ thoả $\mathbf{1}_K^\top v = 0$,
\[
    B v = \alpha v + \beta \mathbf{1}_K (\mathbf{1}_K^\top v) = \alpha v.
\]
Không gian $\{v : \mathbf{1}_K^\top v = 0\}$ có chiều $K-1$ nên $\alpha$ là giá trị riêng bội $K-1$.

Để chuyển sang $CBC^\top$, dùng tính cân bằng $C^\top C = (n/K) I_K$: với mọi vector riêng $v$ của $B$ có giá trị riêng $\mu$,
\[
    (C B C^\top)(C v) = C B (C^\top C) v = \frac{n}{K} \mu \, (C v).
\]
Mặt khác $CBC^\top$ có hạng tối đa $K$ vì chứa thừa số $C \in \mathbb{R}^{n \times K}$, nên các giá trị riêng khác không của $CBC^\top$ chính là $(n/K) \mu$ với $\mu$ chạy trên phổ của $B$.

Thay $\mu = \alpha + K\beta$ và $\mu = \alpha$ cho hai giá trị riêng ở phát biểu.
\end{proof}

\subsection*{Chứng minh Bổ đề~\ref{lem:gap-AE} (Spectral gap của $\mathbb{E}[A_{E^2}]$)}

\begin{proof}
Theo định nghĩa mô hình, $A_{E^2}[u, v] = E_{uv}$ là biến Bernoulli với tham số
\[
    \pi^e_{uv} = \begin{cases}
        a_e/n & \text{nếu } \bar e_u = \bar e_v, \\[2pt]
        b_e/n & \text{nếu } \bar e_u \neq \bar e_v.
    \end{cases}
\]
Đặt $B_{E^2} \in \mathbb{R}^{K \times K}$ với
\[
    B_{E^2}[k, \ell] = \begin{cases}
        a_e/n & \text{nếu } k = \ell, \\[2pt]
        b_e/n & \text{nếu } k \neq \ell.
    \end{cases}
\]
Khi đó
\begin{align*}
    \mathbb{E}[A_{E^2}]
        &= C B_{E^2} C^\top - \frac{a_e}{n} I_n \\
        &\approx C B_{E^2} C^\top,
\end{align*}
trong đó dòng cuối bỏ đường chéo bậc $O(1/n)$, không ảnh hưởng đến hai giá trị riêng chính. Viết $B_{E^2}$ về dạng hai mức
\[
    B_{E^2} = \alpha I_K + \beta \mathbf{1}_K \mathbf{1}_K^\top, \qquad \alpha = \dfrac{a_e - b_e}{n}, \quad \beta = \dfrac{b_e}{n},
\]
rồi áp dụng Bổ đề~\ref{lem:spectrum-CBCt} với cặp $(\alpha, \beta)$ trên, hai giá trị riêng khác không của $\mathbb{E}[A_{E^2}]$ là
\begin{align*}
    \frac{n}{K} (\alpha + K \beta) &= \frac{a_e - b_e}{K} + b_e &&= \frac{a_e + (K-1) b_e}{K} && (\text{bội } 1), \\
    \frac{n}{K} \alpha             &                            &&= \frac{a_e - b_e}{K}        && (\text{bội } K-1).
\end{align*}
Giả thiết $a_e > b_e$ cho cả hai giá trị riêng dương và giá trị nhỏ hơn cho
\[
    \lambda_{\min}^+(\mathbb{E}[A_{E^2}]) = \dfrac{a_e - b_e}{K}.
\]
\end{proof}

\subsection*{Chứng minh Bổ đề~\ref{lem:gap-AET} (Spectral gap của $\mathbb{E}[A_{ET}]$)}

\begin{proof}
Phân rã $A_{ET} = A_{T^2} - D$ với $D := A_{T^2} - A_{ET}$. Vì $A_{T^2}[u, v] = \sum_{w \neq u, v} T_{uvw}$ là tổng các chỉ thị Bernoulli độc lập, tính tuyến tính của kỳ vọng cho ngay
\[
    \mathbb{E}[A_{T^2}[u, v]] = \sum_{w \neq u, v} \pi^t_{uvw} =: \mu_{uv}.
\]
Giá trị $\pi^t_{uvw}$ phụ thuộc cộng đồng của bộ ba $\{u, v, w\}$ nên ta tính $\mu_{uv}$ qua tách tổng thành hai trường hợp theo cộng đồng của cặp $\{u, v\}$.

\emph{Trường hợp 1: $\bar e_u = \bar e_v = k$.} Tách tổng theo việc $w$ có cùng cộng đồng với $\{u, v\}$ hay không:
\begin{align}
    \mu_{uv}
        &= \sum_{w \in C_k \setminus \{u, v\}} \pi^t_{uvw} + \sum_{w \notin C_k} \pi^t_{uvw} \notag \\
        &= \left(\frac{n}{K} - 2\right) \cdot \frac{a_t}{n^2} + \left(n - \frac{n}{K}\right) \cdot \frac{b_t}{n^2} \notag \\
        &= \frac{1}{n} \cdot \frac{a_t + (K-1) b_t}{K} + O(1/n^2). \label{eq:AT2-case-same}
\end{align}
Tổng thứ nhất ứng với bộ ba cùng cộng đồng nên $\pi^t = a_t/n^2$ và có $n/K - 2$ đỉnh $w$ trong $C_k \setminus \{u, v\}$. Tổng thứ hai ứng với bộ ba có $w$ khác cụm nên $\pi^t = b_t/n^2$ và có $n - n/K$ đỉnh $w$ ngoài $C_k$.

\emph{Trường hợp 2: $\bar e_u \neq \bar e_v$.} Đặt $k_1 = \bar e_u$, $k_2 = \bar e_v$. Tách tổng theo việc $w$ thuộc $C_{k_1} \cup C_{k_2}$ hay không:
\begin{align}
    \mu_{uv}
        &= \sum_{w \in (C_{k_1} \cup C_{k_2}) \setminus \{u, v\}} \pi^t_{uvw} + \sum_{w \notin C_{k_1} \cup C_{k_2}} \pi^t_{uvw} \notag \\
        &= \left(\frac{2 n}{K} - 2\right) \cdot \frac{b_t}{n^2} + \left(n - \frac{2 n}{K}\right) \cdot \frac{b_t}{n^2} \notag \\
        &= (n - 2) \cdot \frac{b_t}{n^2}
        = \frac{b_t}{n} + O(1/n^2). \label{eq:AT2-case-diff}
\end{align}
Ở cả hai nhóm, bộ ba $\{u, v, w\}$ đều không cùng cộng đồng (vì đã có $\bar e_u \neq \bar e_v$) nên $\pi^t_{uvw} = b_t/n^2$ đồng nhất.

Gộp hai trường hợp~\eqref{eq:AT2-case-same}, \eqref{eq:AT2-case-diff},
\[
    \mu_{uv} = \begin{cases}
        \dfrac{1}{n} \cdot \dfrac{a_t + (K-1) b_t}{K} + O(1/n^2) & \text{nếu } \bar e_u = \bar e_v, \\[8pt]
        \dfrac{b_t}{n} + O(1/n^2) & \text{nếu } \bar e_u \neq \bar e_v.
    \end{cases}
\]

Tiếp theo tính $\mathbb{E}[D[u, v]]$, ứng với trùng lặp tam giác. Phần tử $D[u, v]$ chỉ khác $0$ khi có từ hai tam giác $G_T$ trở lên qua $(u, v)$. Vì $D = A_{T^2} - A_{ET}$,
\[
    \mathbb{E}[D[u, v]] = \mathbb{E}[A_{T^2}[u, v]] - \mathbb{E}[A_{ET}[u, v]] = \mu_{uv} - \left[1 - \prod_{w \neq u, v} (1 - \pi^t_{uvw})\right].
\]
Khai triển tích $\prod_w (1 - \pi^t_{uvw})$ ta được:
\[
    1 - \prod_{w \neq u, v} (1 - \pi^t_{uvw}) = \sum_w \pi^t_{uvw} - \sum_{w < w'} \pi^t_{uvw} \pi^t_{uvw'} + \sum_{w < w' < w''} \pi^t_{uvw} \pi^t_{uvw'} \pi^t_{uvw''} - \cdots
\]
Đại lượng đầu là $\mu_{uv}$. Tổng các bậc ba trở lên bị chặn bởi
\[
    \sum_{w < w' < w''} \pi^t_{uvw} \pi^t_{uvw'} \pi^t_{uvw''} \leq \frac{1}{6}\Big(\sum_w \pi^t_{uvw}\Big)^3 = \frac{\mu_{uv}^3}{6} = O(\mu_{uv}^3).
\]
Thay vào,
\[
    \mathbb{E}[D[u, v]] = \sum_{w < w'} \pi^t_{uvw} \pi^t_{uvw'} + O(\mu_{uv}^3).
\]
Khai triển bình phương của tổng,
\[
    \Big(\sum_w \pi^t_{uvw}\Big)^2 = \sum_w (\pi^t_{uvw})^2 + 2 \sum_{w < w'} \pi^t_{uvw} \pi^t_{uvw'},
\]
do đó
\[
    \sum_{w < w'} \pi^t_{uvw} \pi^t_{uvw'} = \frac{1}{2}\left[\Big(\sum_w \pi^t_{uvw}\Big)^2 - \sum_w (\pi^t_{uvw})^2\right] = \frac{1}{2} \mu_{uv}^2 - \frac{1}{2} \sum_w (\pi^t_{uvw})^2.
\]
kết hợp
\[
    \sum_w (\pi^t_{uvw})^2 = O(n \cdot 1/n^4) = O(1/n^3) = o(\mu_{uv}^2),
\]
ta được
\[
    \mathbb{E}[D[u, v]] = \frac{1 + o(1)}{2} \mu_{uv}^2 = O(1/n^2),
\]
trong đó dùng $\mu_{uv} = O(1/n)$.

Gộp lại, $\mathbb{E}[A_{ET}[u, v]] = \mathbb{E}[A_{T^2}[u, v]] - \mathbb{E}[D[u, v]] = \mu_{uv} + O(1/n^2)$, tức
\[
    \mathbb{E}[A_{ET}[u, v]] = \begin{cases}
        \dfrac{1}{n} \cdot \dfrac{a_t + (K-1) b_t}{K} + O(1/n^2) & \text{nếu } \bar e_u = \bar e_v, \\[8pt]
        \dfrac{b_t}{n} + O(1/n^2) & \text{nếu } \bar e_u \neq \bar e_v.
    \end{cases}
\]

Dùng lại ma trận chỉ thị $C$ và đặt $B_{ET} \in \mathbb{R}^{K \times K}$ với
\[
    B_{ET}[k, \ell] = \begin{cases}
        \dfrac{a_t + (K-1) b_t}{K n} & \text{nếu } k = \ell, \\[8pt]
        \dfrac{b_t}{n} & \text{nếu } k \neq \ell.
    \end{cases}
\]
Khi đó $\mathbb{E}[A_{ET}] \approx C B_{ET} C^\top$ ở bậc dẫn đầu, sai số chuẩn phổ $O(1/n)$ không ảnh hưởng $\lambda_{\min}^+$. Viết $B_{ET}$ về dạng hai mức
\[
    B_{ET} = \alpha_T I_K + \beta_T \mathbf{1}_K \mathbf{1}_K^\top, \qquad \alpha_T = \dfrac{a_t - b_t}{K n}, \quad \beta_T = \dfrac{b_t}{n},
\]
rồi áp dụng Bổ đề~\ref{lem:spectrum-CBCt} cho hai giá trị riêng khác không là $\tfrac{a_t + (K^2 - 1) b_t}{K^2}$ (bội $1$) và $\tfrac{a_t - b_t}{K^2}$ (bội $K-1$). Giả thiết $a_t > b_t$ cho cả hai dương, và giá trị nhỏ hơn cho gap.
\end{proof}

\subsection*{Chứng minh Bổ đề~\ref{lem:bound-AET} (Bound chuẩn phổ của $A_{ET}$)}

\begin{proof}
$A_{ET}$ là ma trận kề của các cạnh quan sát được trong $G_T$. Cùng một chỉ thị $T_{uvw}$ ảnh hưởng đến cả ba phần tử $A_{ET}[u, v]$, $A_{ET}[u, w]$, $A_{ET}[v, w]$, nên các phần tử của $A_{ET}$ không độc lập theo cặp $(u, v)$. Để tận dụng đánh giá của Paul, Milenkovic và Chen \cite{paul2023hosc} cho ma trận đếm $A_{T^2}$, ta phân rã
\[
    A_{ET} = A_{T^2} - D, \qquad D = A_{T^2} - A_{ET}.
\]
Bất đẳng thức tam giác cho
\[
    \|A_{ET} - \mathbb{E}[A_{ET}]\|_2
        \leq \|A_{T^2} - \mathbb{E}[A_{T^2}]\|_2 + \|D - \mathbb{E}[D]\|_2.
\]

Đại lượng đầu có bound
\[
    \|A_{T^2} - \mathbb{E}[A_{T^2}]\|_2 \lesssim \sqrt{a_t}
\]
với xác suất $1 - o(1)$, trích trực tiếp từ kết quả của Paul, Milenkovic và Chen \cite{paul2023hosc} cho $A_{T^2}$.

Còn lại $\|D - \mathbb{E}[D]\|_2$. Dùng $\|M\|_2 \leq \|M\|_F$ cho ma trận tuỳ ý,
\[
    \|D - \mathbb{E}[D]\|_2 \leq \|D - \mathbb{E}[D]\|_F.
\]
Lấy kỳ vọng bình phương và đổi tổng,
\[
    \mathbb{E}\big[\|D - \mathbb{E}[D]\|_F^2\big] = \sum_{u, v} \mathrm{Var}(D[u, v]).
\]
Vậy chỉ cần chặn phương sai từng phần tử.

Đặt $X_{uv} := A_{T^2}[u, v] = \sum_w T_{uvw}$. Theo định nghĩa,
\[
    D[u, v] = \begin{cases}
        X_{uv} - 1 & \text{nếu } X_{uv} \geq 2, \\
        0 & \text{nếu trái lại}.
    \end{cases}
\]
Khi $X_{uv} \geq 2$, $(X_{uv} - 1)^2 \leq X_{uv}(X_{uv} - 1)$, do đó $D[u, v]^2 \leq X_{uv}(X_{uv} - 1)$ ở mọi giá trị của $X_{uv}$. Vì các $T_{uvw}$ độc lập theo $w$,
\[
    \mathbb{E}[X_{uv}(X_{uv} - 1)] = \sum_{w \neq w'} \pi^t_{uvw} \pi^t_{uvw'} \leq \Big(\sum_w \pi^t_{uvw}\Big)^2 = \mu_{uv}^2.
\]
Kết hợp,
\[
    \mathrm{Var}(D[u, v]) \leq \mathbb{E}[D[u, v]^2] \leq \mu_{uv}^2 = O(1/n^2).
\]

Tổng theo $n^2$ cặp $(u, v)$,
\[
    \mathbb{E}\big[\|D - \mathbb{E}[D]\|_F^2\big] = \sum_{u, v} \mathrm{Var}(D[u, v]) \leq n^2 \cdot O(1/n^2) = O(1).
\]
Đặt $Y := \|D - \mathbb{E}[D]\|_F^2$. Áp dụng bất đẳng thức Markov cho biến ngẫu nhiên không âm $Y$: với mỗi $M > 0$,
\[
    \Pr(Y > M) \leq \frac{\mathbb{E}[Y]}{M} = \frac{O(1)}{M}.
\]
Chọn $M = \log n$, xác suất ở vế phải tiến đến $0$. Do đó với xác suất $1 - o(1)$,
\[
    \|D - \mathbb{E}[D]\|_F = O(\sqrt{\log n}) = o(\sqrt{a_t}),
\]
trong đó dùng giả thiết $a_t \gtrsim \log n$. Dùng tiếp $\|M\|_2 \leq \|M\|_F$,
\[
    \|D - \mathbb{E}[D]\|_2 = o(\sqrt{a_t}).
\]

Kết hợp với $\|A_{T^2} - \mathbb{E}[A_{T^2}]\|_2 \lesssim \sqrt{a_t}$,
\[
    \|A_{ET} - \mathbb{E}[A_{ET}]\|_2 \leq \|A_{T^2} - \mathbb{E}[A_{T^2}]\|_2 + \|D - \mathbb{E}[D]\|_2 \lesssim \sqrt{a_t}
\]
với xác suất $1 - o(1)$.
\end{proof}

\subsection*{Chứng minh Mệnh đề~\ref{prop:lambda-opt} (Tham số nội suy tối ưu)}

\begin{proof}
Đặt
\[
    N(\lambda) := a_e + \lambda^2 \tilde a_t, \qquad D(\lambda) := \alpha + \beta + \lambda \tilde\beta,
\]
nên $f(\lambda) = N(\lambda) / D(\lambda)^2$. Đạo hàm
\[
    N'(\lambda) = 2 \lambda \tilde a_t, \qquad D'(\lambda) = \tilde\beta.
\]
Quy tắc thương cho
\[
    f'(\lambda) = \frac{N'(\lambda) D(\lambda) - 2 N(\lambda) D'(\lambda)}{D(\lambda)^3}.
\]
Vì $D(\lambda) > 0$ với mọi $\lambda \geq 0$, điểm dừng $f'(\lambda) = 0$ tương đương tử số bằng $0$:
\[
    N'(\lambda) D(\lambda) = 2 N(\lambda) D'(\lambda).
\]
Thay các biểu thức của $N, N', D, D'$ và chia hai vế cho $2$, đẳng thức tương đương với
\begin{align*}
    \lambda \tilde a_t \left[\alpha + \beta + \lambda \tilde\beta\right] &= (a_e + \lambda^2 \tilde a_t) \tilde\beta \\
    \iff \lambda \tilde a_t (\alpha + \beta) + \lambda^2 \tilde a_t \tilde\beta &= a_e \tilde\beta + \lambda^2 \tilde a_t \tilde\beta \\
    \iff \lambda \tilde a_t (\alpha + \beta) &= a_e \tilde\beta \\
    \iff \lambda &= \frac{a_e \tilde\beta}{\tilde a_t (\alpha + \beta)},
\end{align*}
trong đó dòng thứ hai triệt tiêu $\lambda^2 \tilde a_t \tilde\beta$ ở hai vế. Đây là điểm $\lambda^*$ ở phát biểu. Kiểm tra giá trị biên $f(0) = a_e/(\alpha + \beta)^2$ và $f(\lambda) \to \tilde a_t/\tilde\beta^2$ khi $\lambda \to \infty$ cho thấy đây là cực tiểu duy nhất trên $\lambda \geq 0$.
\end{proof}

